\begin{figure}[htbp]
\centering
\begin{tikzpicture}[scale=0.9]
% Define colors for points
\definecolor{pointA}{RGB}{255, 0, 0}      % Red
\definecolor{pointB}{RGB}{0, 150, 0}      % Green  
\definecolor{pointC}{RGB}{0, 0, 255}      % Blue
\definecolor{pointD}{RGB}{255, 150, 0}    % Orange

% Frame 1 (t)
\node[anchor=south] at (2, 5) {\textbf{Frame at time t}};
\draw[dashed, thick] (-0.5, -0.5) rectangle (4.5, 4.5);
\fill[white] (-0.5, -0.5) rectangle (4.5, 4.5);

% Points in frame 1
\fill[pointA] (1,3) circle (0.15);     % Point A
\fill[pointB] (3,3) circle (0.15);     % Point B  
\fill[pointC] (1,1) circle (0.15);     % Point C
\fill[pointD] (3,1) circle (0.15);     % Point D

% Labels for points in frame 1
\node[above] at (1,3.2) {\textbf{a}};
\node[above] at (3,3.2) {\textbf{b}};
\node[below] at (1,0.8) {\textbf{c}};
\node[below] at (3,0.8) {\textbf{d}};

% Motion vectors from frame 1
\draw[-{Latex[length=3mm]}, thick, black] (1,3) -- (1.8,2.2);     % A moves right-down
\draw[-{Latex[length=3mm]}, thick, black] (3,3) -- (2.2,3.8);     % B moves left-up  
\draw[-{Latex[length=3mm]}, thick, black] (1,1) -- (0.2,1.8);     % C moves left-up
\draw[-{Latex[length=3mm]}, thick, black] (3,1) -- (3.8,0.2);     % D moves right-down

% Frame 2 (t + δt)
\node[anchor=south] at (8, 5) {\textbf{Frame at time t + $\delta$t}};
\draw[dashed, thick] (5.5, -0.5) rectangle (10.5, 4.5);
\fill[white] (5.5, -0.5) rectangle (10.5, 4.5);

% Points in frame 2 - same colors (brightness constancy)
\fill[pointA] (7.8,2.2) circle (0.15);   % Point A moved
\fill[pointB] (8.2,3.8) circle (0.15);   % Point B moved
\fill[pointC] (6.2,1.8) circle (0.15);   % Point C moved  
\fill[pointD] (9.8,0.2) circle (0.15);   % Point D moved

% Labels for points in frame 2
\node[above right] at (7.8,2.35) {\textbf{a}};
\node[above left] at (8.2,3.95) {\textbf{b}};
\node[above left] at (6.2,1.95) {\textbf{c}};
\node[below left] at (9.8,0.05) {\textbf{d}};

% Brightness constancy equations
\node[anchor=north, align=center] at (5.0, -0.5) {
    $I(\mathbf{a}(t),t) = I(\mathbf{a}(t+\delta t),t+\delta t)$\\
    $I(\mathbf{b}(t),t) = I(\mathbf{b}(t+\delta t),t+\delta t)$\\
    $I(\mathbf{c}(t),t) = I(\mathbf{c}(t+\delta t),t+\delta t)$\\
    $I(\mathbf{d}(t),t) = I(\mathbf{d}(t+\delta t),t+\delta t)$
};

% Arrow between frames
\draw[-{Latex[length=2mm]}, thick] (4.6,2) -- (5.4,2);
\node[above] at (5.0,2.2) {$\delta t$};

\end{tikzpicture}
\caption{Diagram representing the optical flow problem. Coloured points represent piexls. $I$: brightness/intensity of a pixel at time $t$. Adapted from Gkioulekas~\cite{gkioulekas2019}.}
\label{fig:of_diagram}
\end{figure}